\appendix

\section{Full Reduction Rules}
\label{app:rules}

This appendix gives the reduction rules elided from \sect{sec:sem-reduction}. Evaluation contexts $E$ are the standard left-to-right contexts over sequencing, conditionals, loops, assignment right-hand sides, and primitive arguments.

\subsection{Pure Core (L1)}
\label{app:rules-core}

\begin{mathpar}
\inferrule*[right=\rulename{Assign}]
  {\MEnv(x) = (\_,\tau,p,\texttt{var})}
  {\MConfig{\MEnv}{\MErr}{\MPol}{\MPrin}{\Assign{x}{v}} \step \MConfig{\MEnv[\MBind{x}{v}{\tau}{p}{\texttt{var}}]}{\MErr}{\MPol}{\MPrin}{\Tunit}}

\inferrule*[right=\rulename{Var-Bind}]
  { }
  {\MConfig{\MEnv}{\MErr}{\MPol}{\MPrin}{\Var{x}{\tau}{v}} \step \MConfig{\MEnv[\MBind{x}{v}{\tau}{\emptyset}{\texttt{var}}]}{\emptyset}{\MPol}{\MPrin}{\Tunit}}

\inferrule*[right=\rulename{If-True}]
  { }
  {\MConfig{\MEnv}{\MErr}{\MPol}{\MPrin}{\Iff{\texttt{true}}{e_1}{e_2}} \step \MConfig{\MEnv}{\MErr}{\MPol}{\MPrin}{e_1}}

\inferrule*[right=\rulename{If-False}]
  { }
  {\MConfig{\MEnv}{\MErr}{\MPol}{\MPrin}{\Iff{\texttt{false}}{e_1}{e_2}} \step \MConfig{\MEnv}{\MErr}{\MPol}{\MPrin}{e_2}}

\inferrule*[right=\rulename{While}]
  { }
  {\MConfig{\MEnv}{\MErr}{\MPol}{\MPrin}{\Whl{e}{e'}} \step \MConfig{\MEnv}{\MErr}{\MPol}{\MPrin}{\Iff{e}{\Seq{e'}{\Whl{e}{e'}}}{\Tunit}}}

\inferrule*[right=\rulename{For-Nil}]
  { }
  {\MConfig{\MEnv}{\MErr}{\MPol}{\MPrin}{\Forr{x}{[]}{e}} \step \MConfig{\MEnv}{\MErr}{\MPol}{\MPrin}{\Tunit}}

\inferrule*[right=\rulename{For-Cons}]
  { }
  {\MConfig{\MEnv}{\MErr}{\MPol}{\MPrin}{\Forr{x}{v::vs}{e}} \step \MConfig{\MEnv[\MBind{x}{v}{\tau}{\emptyset}{\texttt{let}}]}{\MErr}{\MPol}{\MPrin}{\Seq{e}{\Forr{x}{vs}{e}}}}

\inferrule*[right=\rulename{Seq}]
  { }
  {\MConfig{\MEnv}{\MErr}{\MPol}{\MPrin}{\Seq{v}{e}} \step \MConfig{\MEnv}{\MErr}{\MPol}{\MPrin}{e}}

\inferrule*[right=\rulename{Js}]
  {\JOracle{\MEnv}{c} = v}
  {\MConfig{\MEnv}{\MErr}{\MPol}{\MPrin}{\Js{c}} \step \MConfig{\MEnv}{\MErr}{\MPol}{\MPrin}{v}}
\end{mathpar}

\subsection{Ask, Say, Agent, Eval (L1/L2)}
\label{app:rules-io}

\begin{mathpar}
\inferrule*[right=\rulename{Ask}]
  {\text{human provides }v : \Tstring}
  {\MConfig{\MEnv}{\MErr}{\MPol}{\MPrin}{\Ask{s}} \step \MConfig{\MEnv}{\MErr}{\MPol}{\MPrin}{v}}

\inferrule*[right=\rulename{Say}]
  {\text{emit }s\text{ to output stream}}
  {\MConfig{\MEnv}{\MErr}{\MPol}{\MPrin}{\Say{s}} \step \MConfig{\MEnv}{\MErr}{\MPol}{\MPrin}{\Tunit}}

\inferrule*[right=\rulename{Agent-Success}]
  {\PolOk{\MPol}{\MPrin}{\texttt{agent}} \\
   \Oracle{s}{\MErr}{\Tstring} \ni v \\
   \RtType{\MEnv}{v}{\Tstring}}
  {\MConfig{\MEnv}{\MErr}{\MPol}{\MPrin}{E[\Agent{s}]} \step \MConfig{\MEnv}{\emptyset}{\MPol}{\MPrin}{E[v]}}

\inferrule*[right=\rulename{Eval-Success}]
  {\PolOk{\MPol}{\MPrin}{\texttt{eval}} \\
   v = \MEnv(f) \\
   v = \texttt{workflow}\;f(\overline{x_i{:}T_i}) \to T\;\{e_b\} \\
   \forall i.\,\RtType{\MEnv}{v_i}{T_i}}
  {\MConfig{\MEnv}{\MErr}{\MPol}{\MPrin}{E[\Eval{f}{\overline{v_i}}]} \step \MConfig{\MEnv[\overline{\MBind{x_i}{v_i}{T_i}{\emptyset}{\texttt{let}}}]}{\MErr}{\MPol}{\MPrin}{E[e_b]}}
\end{mathpar}

The $\rulename{Agent-Success}$ rule mirrors $\rulename{Gen-Success}$ but always returns $\Tstring$, so no residual re-type-check is needed. $\rulename{Eval-Success}$ invokes a workflow value that was either defined statically or produced by a prior arrow-typed $\Gen{\Tarrow{A}{R}}{s}$; argument types are checked at the call site and the body $e_b$ reduces in the caller's environment extended with the parameter bindings.

\subsection{Agent and Eval Heal}
\label{app:rules-heal}

Heal rules for $\Agent{s}$ and $\Eval{f}{\overline{v_i}}$ mirror $\rulename{Gen-Heal-Type}$ and $\rulename{Gen-Heal-Pol}$:

\begin{mathpar}
\inferrule*[right=\rulename{Agent-Heal-Type}]
  {\Oracle{s}{\MErr}{\Tstring} \ni v \\
   \lnot\,\RtType{\MEnv}{v}{\Tstring} \\
   \epsilon = \Explain{v}{\Tstring} \\
   |\MErr \cdot \epsilon| \le \MRetry}
  {\MConfig{\MEnv}{\MErr}{\MPol}{\MPrin}{E[\Agent{s}]} \step \MConfig{\MEnv}{\MErr \cdot \epsilon}{\MPol}{\MPrin}{E[\Agent{s}]}}

\inferrule*[right=\rulename{Enforce-Heal}]
  {\MEnv(x) = (p,\Tpolicy,\_,\_) \\
   \Install{\MPol}{p} \text{ undefined or widens} \\
   \epsilon = \Explain{p}{\MPol} \\
   |\MErr \cdot \epsilon| \le \MRetry}
  {\MConfig{\MEnv}{\MErr}{\MPol}{\MPrin}{E[\Enforce{x}]} \step \MConfig{\MEnv}{\MErr \cdot \epsilon}{\MPol}{\MPrin}{E[\Gen{\Tpolicy}{s_x}]}}
\end{mathpar}

In $\rulename{Enforce-Heal}$, $s_x$ denotes the prompt recorded in the provenance of the binding for $x$. The rule rewinds to the causal gen so the oracle can retry with the Cedar compilation or monotonicity error as feedback.

\section{Proof of Theorem~\ref{thm:hadl-sound}}
\label{app:proofs}

We expand the argument clause by clause. Each induction is on the derivation of $C_n \step C_{n+1}$, mirroring the Lean case split on the constructor of $\mathtt{Step}$. The mechanization relies on five axioms: \texttt{jsEval} and \texttt{jsEval\_wellTyped} (the JavaScript host evaluates pure expressions and returns runtime values of the declared type), \texttt{policyInstall\_shrinks} (Cedar's \texttt{forbid}-over-\texttt{permit} semantics implies appended rules shrink $\Allow{\cdot}$), \texttt{freshLabel\_is\_fresh} (the opaque label generator returns names absent from the current environment), and \texttt{freshLabel\_is\_reserved} (generated labels live in the reserved \texttt{\_\_ex\_} namespace; see \S\ref{app:extract}); all other steps are sorry-free.

\subsection{(T1) WF-Preservation}

\paragraph{Case $\rulename{Gen-Success}$.} Assume $\WF{C_n}$. We must show $\WF{C_{n+1}}$. Four side conditions of the rule discharge the four clauses of well-formedness:
\begin{itemize}
  \item $\RtType{\MEnv_n}{v}{\tau}$ establishes clause (1) for the new binding.
  \item $\StType{\Projenv{\MEnv_{n+1}}}{E[v]}{\tau'}$ is exactly clause (2).
  \item The provenance label $\ell$ for the new binding is fresh (axiom \texttt{freshLabel\_is\_fresh}) and cannot create a cycle, establishing clause (3).
  \item $\MErr_{n+1} = \emptyset$, trivially within budget, establishing clause (4).
\end{itemize}

\paragraph{Case $\rulename{Gen-Heal-Type}$, $\rulename{Gen-Heal-Pol}$, $\rulename{Agent-Heal-*}$, $\rulename{Eval-Heal-*}$, $\rulename{Enforce-Heal}$.} $\MEnv$ is unchanged, so clauses (1)--(3) carry over; clause (4) follows from the rule's own $|\MErr \cdot \epsilon| \le \MRetry$ side condition.

\paragraph{Case $\rulename{Gen-Budget-Exhausted}$.} The target configuration has expression $\mathsf{errTerm}(\MErr,\ell)$ and is a terminal-error state; by the statement of (T1) we need not prove $\WF{\cdot}$ for terminal errors.

\paragraph{Case $\rulename{Let-Bind}$ (and $\rulename{Var-Bind}$).} The new binding $\MBind{x}{v}{\tau}{\ell}{\_}$ is well-typed by hypothesis on the evaluation of the right-hand side (from which the rule is only applicable); the residual expression after substitution type-checks because the surrounding let was well-typed at $\tau$.

\paragraph{Case $\rulename{Assign}$.} By the rule, $\MEnv(x) = (\_,\tau,\_,\texttt{var})$ and the assignment writes value $v$ at the same type $\tau$. Clauses (1)--(3) are preserved by the targeted-update lemma (Lean: \texttt{Env.assign\_type\_preserved}); $\MErr$ is unchanged.

\paragraph{Pure-core cases ($\rulename{Var}$, $\rulename{If-*}$, $\rulename{While}$, $\rulename{For-*}$, $\rulename{Seq}$, $\rulename{Js}$).} The environment is unchanged or extended by an iteration binding at the inferred type; $\MPol$ and $\MPrin$ are unchanged. Clauses (1), (3), (4) are trivial. Clause (2) is preserved by the residual typechecker (black-boxed in the mechanization and re-invoked at each $\rulename{Gen-Success}$).

\paragraph{Case $\rulename{Enforce-Install}$.} $\MEnv$ is unchanged; $\MPol$ is updated to $\MPol'$ with $\Allow{\MPol'}\subseteq\Allow{\MPol}$. No binding changes type; provenance and $\MErr$ are unchanged.

\paragraph{Cases $\rulename{Agent-Success}$, $\rulename{Eval-Success}$.} Analogous to $\rulename{Gen-Success}$: each carries a runtime-type check and a residual re-check as side conditions; $\MErr$ is flushed.

\subsection{(T2) Staged Materialization Soundness}

The side condition $\StType{\Projenv{\MEnv_{n+1}}}{E[v]}{\tau'}$ on $\rulename{Gen-Success}$ re-checks the entire residual expression against the extended static context. If the newly bound name $x_\ell$ appears in any subsequent type annotation position, $\Projenv{\MEnv_{n+1}}$ maps $x_\ell$ to the materialized type $\tau$. The $\StType{\cdot}{\cdot}{\cdot}$ judgment resolves name references through this map, so any downstream $\Let{y}{x_\ell}{\Gen{x_\ell}{s'}}$ is checked against the concrete $\tau$ rather than against $\Tschema$.

\subsection{(T3) Policy Monotonicity}

$\MPol$ is updated only by $\rulename{Enforce-Install}$, whose side condition enforces $\Allow{\MPol'} \subseteq \Allow{\MPol}$; this is where the mechanization invokes \texttt{policyInstall\_shrinks}. All other rules leave $\MPol$ unchanged. Monotonicity across the whole trace follows by transitivity. The pointwise authorization claim is immediate because each of $\rulename{Gen-Success}$, $\rulename{Agent-Success}$, and $\rulename{Eval-Success}$ has $\PolOk{\MPol}{\MPrin}{a}$ as a side condition.

\subsection{(T4a) Budget Progress}

If at $C_n$ the redex is $\Gen{\tau}{s}$ and $|\MErr_n|>\MRetry$, no success or heal rule is applicable: the success rules require $\PolOk{\cdot}{\cdot}{\texttt{gen}}\wedge\RtType{\cdot}{\cdot}{\cdot}$ to hold for \emph{some} oracle response, which need not exist; the heal rules have $|\MErr\cdot\epsilon|\le\MRetry$ as a side condition, which fails. Only $\rulename{Gen-Budget-Exhausted}$ applies, and it deterministically transitions to the terminal-error configuration carrying $\MErr_n$ and the fresh provenance label.

\subsection{(T4b) Truthful Gen-Success}

Fix a gen-stuck configuration $C_n$ with $\PolOk{\MPol_n}{\MPrin_n}{\texttt{gen}}$. By eventually-truthfulness (\defn{def:oracle-truthful}) instantiated at $(s,\MEnv_n,\tau,\MPol_n,\MPrin_n)$, there exists $\epsilon$ with $|\epsilon|\le\MRetry$ and $v$ with $\Oracle{s}{\epsilon}{\tau}\ni v$ and $\RtType{\MEnv_n}{v}{\tau}$. The residual static-typing premise is discharged by the $\Tschema$/value wildcards of $\StType{\cdot}{\cdot}{\cdot}$, and freshness of $\ell$ by axiom. All four side conditions of $\rulename{Gen-Success}$ are met, so the rule fires and flushes $\MErr$ to $\emptyset$.

\subsection{(T4c) Oracle-Relative Termination}

Argue by a well-founded measure on reduction sequences. Trace-level gen progress is mechanized (Appendix~\ref{app:extract}): given a hygienic source program, $\mathtt{Extract}$-based defunctionalization lifts (T4b) from a bare redex to any gen-reachable configuration. The remaining ingredient --- a well-founded measure on intermediate pure-core reductions --- is argued on paper. Between two successive applications of a generative rule, only pure-core rules fire. Pure-core reduction terminates on any expression whose \texttt{while} and \texttt{for} loops terminate, which we assume of $e_0$. At each generative step, either a success rule fires or a heal rule extends $\MErr$ by one element. Since $|\MErr| \le \MRetry$, at most $\MRetry$ heal steps occur before either (T4b) fires or $|\MErr|>\MRetry$ forces (T4a). Under eventually-truthful $\mathcal{O}$, (T4b) fires within the budget, $\MErr$ is flushed, and the process repeats for the next generative site. Since the number of generative sites reached along the trace is finite (bounded by the size of $e_0$ plus finitely many eval expansions), the whole trace terminates. When $\mathcal{O}$ is not eventually-truthful at some site, (T4a) fires and produces a terminal error configuration whose provenance label, under the well-formedness assumption of Cor.~\ref{cor:blame}, is exactly the failing gen.

\subsection{Mechanizing $E[\cdot]$ via Extract}
\label{app:extract}

The $E[\cdot]$-congruence cases of \figr{fig:rules-core} and Appendix~\ref{app:rules} are presented on paper as schematic abstractions; the mechanized \texttt{Step} relation is their top-level fragment (all $E=[\cdot]$). To close this gap for trace-level progress --- in particular, to lift (T4b) from ``the redex \emph{is} a gen'' to ``the expression \emph{contains} a gen at an evaluation position'' --- the mechanization \emph{defunctionalizes} $E[\cdot]$ as a syntactic operation.

\paragraph{The \texttt{Extract} operator.} We define a total function
\[
\mathtt{Extract} : \mathsf{Expr} \to (\mathsf{Expr} \times \mathsf{Name} \times \mathsf{Expr})_{\bot}
\]
that, given $e$, either returns $\bot$ when no oracle head is reachable in $e$ at an evaluation position, or returns a triple $(e_h, x, e_k)$ where:
\begin{itemize}
  \item $e_h$ is the \emph{leftmost} oracle head in $e$, i.e.\ a $\Gen{\tau}{s}$ or $\Agent{s}$ occurrence reachable by a standard left-to-right evaluation order;
  \item $x$ is a \emph{fresh reserved} continuation binder, drawn from the namespace of names beginning with the prefix \texttt{\_\_ex\_} (concretely \texttt{\_\_ex\_gen} for gen and \texttt{\_\_ex\_ag} for agent);
  \item $e_k$ is the continuation, i.e.\ the expression obtained from $e$ by replacing the leftmost oracle head with $x$.
\end{itemize}
Loops are handled by a purely syntactic rewrite before recursion: $\Whl{e_c}{p}$ is treated as $\Iff{e_c}{\Seq{p}{\Whl{e_c}{p}}}{\Tunit}$ and $\mathtt{Extract}$ recurses into $e_c$ only. This preserves structural recursion (\texttt{Expr}-sized) and avoids a divergent re-unrolling.

\paragraph{The \texttt{GStep} relation.} On top of the existing pure-core $\step$ relation, we layer $\step_g$:
\[
\inferrule*[Left=\rulename{GStep-Pure}]
  {\mathtt{Extract}(e) = \bot \\ C \step C'}
  {C \step_g C'}
\]
\[
\inferrule*[Left=\rulename{GStep-Run}]
  {\mathtt{Extract}(e) = (e_h, x, e_k) \\
   \MConfig{\MEnv}{\MErr}{\MPol}{\MPrin}{e_h}
     \step
   \MConfig{\MEnv'}{\MErr'}{\MPol'}{\MPrin'}{v} \\
   \RtType{\MEnv'}{v}{\tau} \\
   x \not\in \mathsf{dom}(\MEnv')}
  {\MConfig{\MEnv}{\MErr}{\MPol}{\MPrin}{e}
     \step_g
   \MConfig{\MEnv'[\MBind{x}{v}{\tau}{\bot}{\texttt{let}}]}{\MErr'}{\MPol'}{\MPrin'}{e_k}}
\]
The continuation binder $x$ is \emph{bound by environment extension} rather than substituted. Every downstream use of $x$ inside $e_k$ is resolved by the existing $\rulename{Var}$ rule; no auxiliary substitution machinery is introduced.

\paragraph{Hygiene.} Because continuation binders live in the reserved namespace \texttt{\_\_ex\_*}, the defunctionalization is correct precisely when source programs and oracle-produced values do not name variables in that namespace. The mechanization formalizes this as:
\begin{itemize}
  \item $\mathsf{Name.IsRes}(n) \iff n$ begins with the prefix \texttt{\_\_ex\_};
  \item $\tau.\mathsf{hygienic}$: for a type variable requires that the variable's name is not reserved; trivially true for every other type constructor (the type is inspected only through the $\rulename{tVarResolve}$ rule of $\RtType{\cdot}{\cdot}{\cdot}$);
  \item $\MEnv.\mathsf{Hygienic}$: every $\MEnv$-bound schema's payload type is itself hygienic.
\end{itemize}
The key lemma \texttt{RtType.weaken\_extend\_reserved} states that if $\RtType{\MEnv}{v}{\tau}$ with $\tau$ hygienic and $\MEnv$ hygienic, then extending $\MEnv$ with any reserved-named binding preserves the runtime-type judgment. This is the technical closure for handling the label binding introduced by $\rulename{Gen-Success}$: the new axiom \texttt{freshLabel\_is\_reserved} asserts that generated labels inhabit the same reserved namespace, which is an implementation choice (e.g.\ labels of the form \texttt{\_\_ex\_lbl\_$n$}).

\paragraph{Trace-level gen progress.} With these pieces in place, the mechanization proves
\begin{theorem}[Extract-based gen progress, Lean \texttt{T4\_gStep\_progress\_gen}]
  Let $\mathcal{O}$ be eventually-truthful at budget $\MRetry$ and $\PolOk{\MPol}{\MPrin}{\texttt{gen}}$. If $\mathtt{Extract}(e) = (\Gen{\tau}{s}, x, e_k)$, and the source program is hygienic in the sense above, then there exist $\epsilon$ and $C'$ with $\MConfig{\MEnv}{\epsilon}{\MPol}{\MPrin}{e} \step_g C'$.
\end{theorem}
This is the trace-level lifting of (T4b): no congruence-rule plumbing is needed; the Extract defunctionalization does the routing. The agent-case is structurally identical (same proof shape, different reduction rule, reserved binder \texttt{\_\_ex\_ag}). The pure-core termination measure remaining in (T4c) is orthogonal to this closure and is still argued only on paper.
